% Part: many-valued-logic
% Chapter: three-valued-logics
% Section: goedel

\documentclass[../../../include/open-logic-section]{subfiles}

\begin{document}

\olfileid{mvl}{thr}{god}

\olsection{గోడెల్ తర్కాలు}

కుర్ట్ గోడెల్ $n$-విలువల తర్కాల క్రమాన్ని
ప్రవేశపెట్టాడు. వాటిలో ప్రతి తర్కంలోనూ
అంతర్బోధవాద తర్కంలో చెల్లే అన్ని
\tetoken{సూత్రాలు}{!!{formula}s} ఉంటాయి;
ప్రతి తర్కం సాంప్రదాయిక తర్కంలో ఇమిడి
ఉంటుంది. ఆ క్రమంలో ఆసక్తికరమైన మొదటి తర్కం ఇది:

\begin{defn}\ollabel{defn:goedel}
\emph{$3$-విలువల గోడెల్ తర్కం}~$\LogGod$ను
కింది మాత్రికతో నిర్వచిస్తాం:
\begin{enumerate}
  \item $\lfalse$, $\lnot$, $\land$, $\lor$, $\lif$
  సంయోజకాలు గల ప్రామాణిక ప్రతిజ్ఞావాక్య భాష
  $\Lang L_0$.
  \item సత్యమూల్యాల సమితి $V = \{\True, \Undef, \False\}$.
  \item $\True$ మాత్రమే నిర్దేశిత విలువ; అంటే
  $V^+ = \{\True\}$.
  \item $\lfalse$కు $\tf{\lfalse} = \False$.
  మిగిలిన సంయోజకాల సత్యమూల్య ప్రమేయాలను
  కింది పట్టికలు ఇస్తాయి:
  \begin{center}
    \begin{tabular}{c|c}
      $\tf{\lnot}[\LogGod]$ & \\
      \hline
      $\True$ & $\False$ \\
      $\Undef$ & $\False$ \\
      $\False$ & $\True$
    \end{tabular}
    \quad
    \begin{tabular}{c|ccc}
      $\tf{\land}[\LogGod]$ & $\True$ & $\Undef$ & $\False$ \\
      \hline
      $\True$ & $\True$ & $\Undef$ & $\False$ \\
      $\Undef$ & $\Undef$ & $\Undef$ & $\False$\\
      $\False$ & $\False$ & $\False$ & $\False$
    \end{tabular}
    \\[2ex]
    \begin{tabular}{c|ccc}
      $\tf{\lor}[\LogGod]$ & $\True$ & $\Undef$ & $\False$ \\
      \hline
      $\True$ & $\True$ & $\True$ & $\True$ \\
      $\Undef$ & $\True$ & $\Undef$ & $\Undef$ \\
      $\False$ & $\True$ & $\Undef$ & $\False$
    \end{tabular}
    \quad
    \begin{tabular}{c|ccc}
      $\tf{\lif}[\LogGod]$ & $\True$ & $\Undef$ & $\False$ \\
      \hline
      $\True$ & $\True$ & $\Undef$ & $\False$ \\
      $\Undef$ & $\True$ & $\True$ & $\False$  \\
      $\False$ & $\True$ & $\True$ & $\True$
    \end{tabular}
  \end{center}
\end{enumerate}
\end{defn}

$\land$, $\lor$ల సత్య పట్టికలు \L ukasiewicz
తర్కంలోనూ బలమైన క్లీని తర్కంలోనూ ఉన్నట్లే
ఉన్నాయి; కానీ $\lnot$, $\lif$ల పట్టికలు
ప్రతి తర్కంలో వేరుగా ఉన్నాయి. గోడెల్ తర్కంలో
$\tf{\lnot}(\Undef) = \False$.
\L ukasiewicz, క్లీని తర్కాలకు భిన్నంగా
$\tf{\lif}(\Undef, \False) = \False$;
క్లీని తర్కానికి భిన్నంగా, కానీ
\L ukasiewicz తర్కంలో ఉన్నట్లుగా,
$\tf{\lif}(\Undef,
\Undef) = \True$.

పైన సూచించిన అంతర్బోధవాద తర్కంతో సంబంధం
తెలియజేసినట్లుగా, $\LogGod[3]$
అంతర్బోధవాద తర్కానికి దగ్గరగా ఉంటుంది.
అంతర్బోధవాద తర్కంలోని అన్ని సత్యాలూ
~$\LogGod[3]$లో సర్వసత్యాలు. అంతర్బోధవాద
తర్కంలో చెల్లని అనేక సాంప్రదాయిక సర్వసత్యాలు
~$\LogGod[3]$లోనూ సర్వసత్యాలు కావు.
ఉదాహరణకు కిందివి సర్వసత్యాలు కావు:
\begin{align*}
  & p \lor \lnot p && (p \lif q) \lif (\lnot p \lor q) \\
  & \lnot\lnot p \lif p && \lnot(\lnot p \land \lnot q) \lif (p \lor q) \\
  & ((p \lif q) \lif p) \lif p && \lnot(p \lif q) \lif (p \land \lnot q)
\end{align*}
అయితే $\LogGod[3]$లోని ప్రతి సర్వసత్యమూ
అంతర్బోధవాద తర్కంలో చెల్లదు. ఉదాహరణకు
$\lnot\lnot p \lor \lnot p$, $(p
\lif q) \lor (q \lif p)$.

\begin{prob}
  కిందివి~$\LogGod[3]$లో సర్వసత్యాలని
  సత్య పట్టికలతో చూపండి:
  \begin{align*}
    & \lnot\lnot p \lor \lnot p\\
    & (p \lif q) \lor (q \lif p) \\
    & \lnot(p \land q) \lif (\lnot p \lor \lnot q) \\
    & (p \lif q) \lor (q \lif r) \lor (r \lif s)
  \end{align*}
\end{prob}

\begin{prob}
  కిందివి~$\LogGod[3]$లో సర్వసత్యాలు కావని
  సత్య పట్టికలతో చూపండి:
  \begin{align*}
    & (p \lif q) \lif (\lnot p \lor q) \\
    & \lnot(\lnot p \land \lnot q) \lif (p \lor q) \\
    & ((p \lif q) \lif p) \lif p \\
    & \lnot(p \lif q) \lif (p \land \lnot q)
  \end{align*}
\end{prob}

\begin{prob}
  కింది సంబంధాల్లో ఏవి గోడెల్ తర్కంలో
  నిలుస్తాయి? ప్రతిదానికీ సత్య పట్టిక ఇవ్వండి.
  \begin{enumerate}
    \item $p, p \lif q \Entails q$
    \item $p \lor q, \lnot p \Entails q$
    \item $p \land q \Entails p$
    \item $p \Entails p \land p$
    \item $p \Entails p \lor q$
  \end{enumerate}
\end{prob}

\end{document}
